\documentclass[../../main.tex]{subfiles}

\begin{document}

\section{拓扑空间之间的连续映射}\label{section:ch11-连续映射}

我们通过收敛的序列定义度量空间之间的映射的连续性. 映射 $f$ 在 $x$ 连续, 若序列收敛到 $x$, 则该序列的象收敛到 $f(x)$. 我们接着证明这等价于用开球表示的 $\varepsilon$-$\delta$ 准则. 连续性的概念可用以下自然的方式推广到拓扑空间之间的映射.

\begin{definition}[连续映射]\label{definition:ch11-连续映射}
对于拓扑空间 $(X,\mathcal{T})$ 和 $(Y,\mathcal{S})$, 映射 $f:X\to Y$ 称为在点 $x_{0}$ 连续, 若对任何 $f(x_{0})$ 的邻域 $\mathcal{O}$, 存在 $x_{0}$ 的邻域 $U$ 使得 $f(U)\subseteq\mathcal{O}$. 若它在 $X$ 的每一点连续, 我们称它为连续的.
\end{definition}

\begin{proposition}\label{proposition:ch11-连续原象}
拓扑空间 $X$ 和 $Y$ 之间的映射 $f:X\to Y$ 是连续的当且仅当对 $Y$ 中的任何开集 $\mathcal{O}$, 它在 $f$ 下的原象 $f^{-1}(\mathcal{O})$ 是 $X$ 的开子集.
\end{proposition}

\begin{proof}
首先假定 $f$ 是连续的. 令 $\mathcal{O}$ 在 $Y$ 中是开的. 根据命题 1, 为证明 $f^{-1}(\mathcal{O})$ 是开的仅须证明 $f^{-1}(\mathcal{O})$ 的每个点有一个包含于 $f^{-1}(\mathcal{O})$ 的邻域. 令 $x$ 属于 $f^{-1}(\mathcal{O})$. 则根据 $f$ 在 $x$ 的连续性, 存在 $x$ 的邻域, 它被映射到 $\mathcal{O}$, 因此它包含于 $f^{-1}(\mathcal{O})$. 反之, 若 $f^{-1}$ 将开集映到开集, $f$ 在整个 $X$ 连续.
\end{proof}

对于从拓扑空间 $X$ 到拓扑空间 $Y$ 的连续映射 $f$, 根据子空间拓扑的定义, $f$ 到 $X$ 的子空间的限制也是连续的. 我们把下一个命题的证明留作练习.

\begin{proposition}\label{proposition:ch11-连续复合}
拓扑空间之间的连续映射的复合是连续的.
\end{proposition}

\begin{definition}[拓扑的强弱]\label{definition:ch11-拓扑强弱}
给定集合 $X$ 的两个拓扑 $\mathcal{T}_{1}$ 和 $\mathcal{T}_{2}$, 若 $\mathcal{T}_{2}\subseteq\mathcal{T}_{1}$, 我们说 $\mathcal{T}_{2}$ 比 $\mathcal{T}_{1}$ 弱, 也说 $\mathcal{T}_{1}$ 比 $\mathcal{T}_{2}$ 强.
\end{definition}

给定集合 $X$ 的一个覆盖 $\mathcal{S}$, 理解 $X$ 的使得覆盖是开的拓扑是有用的. 当然, $\mathcal{S}$ 关于 $X$ 上的离散拓扑是开的. 事实上, 存在 $X$ 的最弱拓扑使得该覆盖是开的: 它是以 $\mathcal{S}$ 为子基的唯一拓扑. 我们将以下命题的证明留作练习.

\begin{proposition}\label{proposition:ch11-弱拓扑基}
令 $X$ 为非空集合而 $\mathcal{S}$ 是覆盖 $X$ 的任何子集族. $\mathcal{S}$ 的有限交组成的 $X$ 的子集族是 $X$ 的拓扑 $\mathcal{T}$ 的基. 它是包含 $\mathcal{S}$ 的最弱拓扑: 若 $\mathcal{T}'$ 是 $X$ 的任何其他包含 $\mathcal{S}$ 的拓扑, 则 $\mathcal{T}\subseteq\mathcal{T}'$.
\end{proposition}

\begin{definition}[由映射族诱导的弱拓扑]\label{definition:ch11-弱拓扑}
令 $X$ 为非空集合而 $\mathcal{F}=\{f_{\alpha}:X\to X_{\alpha}\}_{\alpha\in\Lambda}$ 是映射族, 其中每个 $X_{\alpha}$ 是拓扑空间. $X$ 的包含集族
\begin{align*}
\{f_{\alpha}^{-1}(\mathcal{O}_{\alpha})\mid f_{\alpha}\in\mathcal{F},\mathcal{O}_{\alpha}\text{ 在 }X_{\alpha}\text{ 中是开的}\}
\end{align*}
的最弱拓扑称为 $X$ 的由 $\mathcal{F}$ 诱导的弱拓扑.
\end{definition}

\begin{proposition}\label{proposition:ch11-弱拓扑最少}
令 $X$ 为非空集合而 $\mathcal{F}=\{f_{\lambda}:X\to X_{\lambda}\}_{\lambda\in\Lambda}$ 是映射族, 其中每个 $X_{\lambda}$ 是拓扑空间. $X$ 的由 $\mathcal{F}$ 诱导的弱拓扑是 $X$ 上的拓扑, 使得每个映射 $f_{\lambda}:X\to X_{\lambda}$ 是连续的拓扑中具有最少个数的集合.
\end{proposition}

\begin{proof}
根据命题 10, 对 $\Lambda$ 中的每个 $\lambda$, $f_{\lambda}:X\to X_{\lambda}$ 是连续的当且仅当 $X_{\lambda}$ 中的每个开集在 $f_{\lambda}$ 的原象在 $X$ 中是开的.
\end{proof}

\begin{definition}[同胚]\label{definition:ch11-同胚}
从拓扑空间 $X$ 到拓扑空间 $Y$ 的连续映射称为是同胚映射, 若它是一对一的, 将 $X$ 映上 $Y$, 且有从 $Y$ 到 $X$ 的连续的逆 $f^{-1}$.
\end{definition}

显然同胚的逆是同胚且同胚的复合是同胚. 两个拓扑空间 $X$ 与 $Y$ 称为是同胚的, 若它们之间存在同胚映射. 这是拓扑空间之间的等价关系, 即它是自反、对称与传递的. 从拓扑的观点, 两个同胚的拓扑空间是不可区分的, 这是由于根据命题 10, 对于 $X$ 映上 $Y$ 的同胚 $f$, 集合 $E$ 在 $X$ 中是开的当且仅当它的象 $f(E)$ 在 $Y$ 中是开的. 同胚的概念在拓扑空间中所起的作用正如同等距在度量空间中、群同构在群中所起的作用. 但这里要注意. 在下一个例子我们证明, 对实数的 Lebesgue 可测集 $E$, $L^{1}(E)$ 同胚于 $L^{2}(E)$.

\begin{example}[Mazur]\label{example:ch11-mazur}
令 $E$ 为实数的 Lebesgue 可测集. 对 $L^{1}(E)$ 中的 $f$, 定义 $E$ 上的函数 $\Phi(f)$ 为 $\Phi(f)(x)=\operatorname{sgn}(f(x))\cdot|f(x)|^{1/2}$. 则 $\Phi(f)$ 属于 $L^{2}(E)$. 我们把证明对任何两个数 $a$ 和 $b$,
\begin{align*}
\left|\operatorname{sgn}(a)\cdot|a|^{1/2}-\operatorname{sgn}(b)\cdot|b|^{1/2}\right|^{2}\leqslant 2\cdot|a-b|
\end{align*}
留作练习.

因此对 $L^{1}(E)$ 中的所有 $f,g$,
\begin{align*}
\|\Phi(f)-\Phi(g)\|_{2}^{2}\leqslant 2\cdot\|f-g\|_{1}.
\end{align*}
由此得出 $\Phi$ 是将 $L^{1}(E)$ 映到 $L^{2}(E)$ 的连续的一一映射. 它也将 $L^{1}(E)$ 映上 $L^{2}(E)$ 且它的逆 $\Phi^{-1}$ 定义为对 $L^{2}(E)$ 中的 $f$, $\Phi^{-1}(f)(x)=\operatorname{sgn}(f(x))\cdot|f(x)|^{2}$. 用习题 38 得出逆映射 $\Phi^{-1}$ 是从 $L^{2}(E)$ 到 $L^{1}(E)$ 的连续映射. 因此 $L^{1}(E)$ 同胚于 $L^{2}(E)$, 这些空间的每个具有由 $L^{p}$ 范数诱导的拓扑.
\end{example}

\begin{exercise}\label{exercise:ch11-23}
令 $f$ 为从拓扑空间 $X$ 到拓扑空间 $Y$ 的映射而 $\mathcal{S}$ 是 $Y$ 上的拓扑的子基. 证明 $f$ 是连续的当且仅当 $\mathcal{S}$ 的每个集在 $f$ 下的原象在 $X$ 中是开的.
\end{exercise}

\begin{exercise}\label{exercise:ch11-24}
令 $X$ 为拓扑空间.
\begin{enumerate}[(i)]
\item 若 $X$ 具有平凡拓扑, 找出所有 $X$ 到 $\mathbf{R}$ 的连续映射.
\item 若 $X$ 具有离散拓扑, 找出所有 $X$ 到 $\mathbf{R}$ 的连续映射.
\item 若 $X$ 具有离散拓扑, 找出所有 $\mathbf{R}$ 到 $X$ 的连续一一映射.
\item 若 $X$ 具有平凡拓扑, 找出所有 $\mathbf{R}$ 到 $X$ 的连续一一映射.
\end{enumerate}
\end{exercise}

\begin{exercise}\label{exercise:ch11-25}
对于拓扑空间 $X$ 和 $Y$, 令 $f$ 将 $X$ 映到 $Y$. 以下哪些断言与 $f$ 的连续性等价? 证明你的答案.
\begin{enumerate}[(i)]
\item $Y$ 的每个闭子集在 $f$ 下的原象在 $X$ 中是闭的.
\item 若 $\mathcal{O}$ 在 $X$ 中是开的, 则 $f(\mathcal{O})$ 在 $Y$ 中是开的.
\item 若 $F$ 在 $X$ 中是闭的, 则 $f(F)$ 在 $Y$ 中是闭的.
\item 对 $X$ 的每个子集 $A$, $f(\overline{A})\subseteq\overline{f(A)}$.
\end{enumerate}
\end{exercise}

\begin{exercise}\label{exercise:ch11-26}
证明命题 11.
\end{exercise}

\begin{exercise}\label{exercise:ch11-27}
证明命题 12.
\end{exercise}

\begin{exercise}\label{exercise:ch11-28}
定义在拓扑空间 $X$ 上的两个实值连续函数的和与积是连续的.
\end{exercise}

\begin{exercise}\label{exercise:ch11-29}
令 $\mathcal{F}$ 为集合 $X$ 上的实值函数族. 找出 $\mathcal{F}$ 上使得 $X$ 作为具有由 $\mathcal{F}$ 诱导的弱拓扑的拓扑空间是 Tychonoff 的必要与充分条件.
\end{exercise}

\begin{exercise}\label{exercise:ch11-30}
对于拓扑空间 $X$ 和 $Y$, 令 $f:X\to Y$ 为从 $X$ 到 $Y$ 的一对一与映上的映射. 证明以下断言等价:
\begin{enumerate}[(i)]
\item $f$ 是 $X$ 映上 $Y$ 的同胚.
\item $X$ 的子集 $E$ 在 $X$ 中是开的当且仅当 $f(E)$ 在 $Y$ 中是开的.
\item $X$ 的子集 $E$ 在 $X$ 中是闭的当且仅当 $f(E)$ 在 $Y$ 中是闭的.
\item 集合的闭包的象是象的闭包, 即对 $X$ 的每个子集 $A$, $f(\overline{A})=\overline{f(A)}$.
\end{enumerate}
\end{exercise}

\begin{exercise}\label{exercise:ch11-31}
对于拓扑空间 $X$ 和 $Y$, 令 $f:X\to Y$ 为从 $X$ 到 $Y$ 的连续映射. 假定 $X$ 是 Hausdorff 的, $Y$ 是 Hausdorff 的吗? 假定 $X$ 是正规的, $Y$ 是正规的吗?
\end{exercise}

\begin{exercise}\label{exercise:ch11-32}
令 $\rho_{1}$ 与 $\rho_{2}$ 为集合 $X$ 上的度量, 分别诱导出拓扑 $\mathcal{T}_{1}$ 与 $\mathcal{T}_{2}$. 若 $\mathcal{T}_{1}=\mathcal{T}_{2}$, 这两个度量必须是等价的吗?
\end{exercise}

\begin{exercise}\label{exercise:ch11-33}
证明同胚的逆是同胚. 两个同胚的复合是同胚.
\end{exercise}

\begin{exercise}\label{exercise:ch11-34}
假定拓扑空间 $X$ 具有性质: $X$ 上的每个连续实值函数取到最小值. 证明任何同胚于 $X$ 的拓扑空间也具有该性质.
\end{exercise}

\begin{exercise}\label{exercise:ch11-35}
假定拓扑空间 $X$ 具有性质: $X$ 上的每个连续实值函数有区间作为它的象. 证明任何同胚于 $X$ 的拓扑空间也具有该性质.
\end{exercise}

\begin{exercise}\label{exercise:ch11-36}
证明 $\mathbf{R}$ 同胚于开有界区间 $(0,1)$, 但不同胚于闭有界区间 $[0,1]$.
\end{exercise}

\begin{exercise}\label{exercise:ch11-37}
令 $X$ 与 $Y$ 为拓扑空间, 考虑从 $X$ 到 $Y$ 的映射 $f$. 假定 $X=X_{1}\cup X_{2}$, 而 $f$ 在拓扑子空间 $X_{1}$ 和 $X_{2}$ 上的限制是连续的. 证明 $f$ 不必在 $X$ 的任何点都连续. 若 $X_{1}$ 与 $X_{2}$ 是开的, 证明 $f$ 在 $X$ 上是连续的. 将此与可测函数以及从限制的可测性继承可测性的情形比较.
\end{exercise}

\begin{exercise}\label{exercise:ch11-38}
证明对任何两个数 $a$ 和 $b$,
\begin{align*}
\left|\operatorname{sgn}(a)\cdot|a|^{2}-\operatorname{sgn}(b)\cdot|b|^{2}\right|\leqslant 2\cdot|a-b|(|a|+|b|).
\end{align*}
\end{exercise}

\end{document}